\section{Related Work}
\label{sec:related}

Multi-agent systems fall at the intersection of game theory, distributed systems, and Artificial Intelligence in general~\citep{shoham_08}, and thus have a rich and diverse literature. Our work builds on and is related to prior work in deep multi-agent reinforcement learning, the centralized training and decentralized execution paradigm, and emergent communication protocols.

\textbf{Multi-Agent Reinforcement Learning (MARL).} Within MARL (see \citet{busoniu_08} for a survey), our work
is related to efforts on using recurrent neural networks to approximate agent policies~\citep{hausknecht_aaai15},
stabilizing algorithms for multi-agent training~\citep{lowe_nips17, foerster_aaai18},
and tasks in novel domains \eg coordination and navigation in 3D environments~\citep{peng_arxiv17,openai_18,jaderberg_arxiv18}.

\textbf{Centralized Training \& Decentralized Execution.} Both~\citet{sukhbaatar_nips16} and~\citet{hoshen_nips17}
adopt a centralized framework at both training and test time -- a central controller processes
local observations from all agents and outputs a probability distribution over joint actions.
In this setting, the controller (\eg a fully-connected network) can be viewed as implicitly
encoding communication. \citet{sukhbaatar_nips16} propose an efficient controller architecture
that is invariant to agent permutations by virtue of weight-sharing and averaging (as in~\citet{zaheer_nips17}),
and can, in principle, also be used in a decentralized manner at test time since each
agent just needs its local state vector and the average of incoming messages to take an action.
Meanwhile,~\citet{hoshen_nips17}
proposes to replace averaging by an attentional mechanism to allow targeted interactions between agents.
While closely related to our communication architecture, this work only considers fully-supervised
one-next-step prediction tasks, while we study the full reinforcement learning problem
with tasks requiring planning over long time horizons.

Moreover, a centralized controller quickly becomes intractable in real-world tasks with many agents and high-dimensional observation spaces \eg navigation in House3D~\citep{house3d}.
To address these weaknesses, we adopt the framework of centralized learning but decentralized execution
(following~\citet{foerster_nips16,lowe_nips17}) and further relax it by allowing agents to communicate.
While agents can use extra information during training, at test time, they pick actions solely based on
local observations and communication messages.


\textbf{Emergent Communication Protocols.} Our work is also related to recent work on learning communication
protocols in a completely end-to-end manner with reinforcement learning -- from perceptual input (\eg pixels)
to communication symbols (discrete or continuous) to actions (\eg navigating in an environment).
While~\citep{foerster_nips16,jorge_iclrw16,visdial_rl,kottur_emnlp17,mordatch_arxiv17,lazaridou_iclr17}
constrain agents to communicate with discrete symbols with the explicit goal to study emergence of language,
our work operates in the paradigm of learning a continuous communication protocol in order to solve a downstream
task~\citep{sukhbaatar_nips16,hoshen_nips17,jiang_nips18,singh_iclr19}.
\citet{jiang_nips18,singh_iclr19} also operate in a decentralized execution
setting and use an attentional communication mechanism, but in contrast to our work,
they use attention to decide \emph{when} to communicate, not \emph{who} to communicate with.
In \Secref{sec:competitive}, we discuss how to potentially combine the two approaches.

Table~\ref{table:related} summarizes
the main axes of comparison between our work and previous efforts in this exciting space.

